\documentclass[10pt]{article}
\usepackage[T1]{fontenc}
\usepackage[utf8]{inputenc}
\usepackage[french]{babel}
\usepackage{amsmath,amssymb,mathrsfs,mathpazo}
\usepackage{geometry}
\geometry{a4paper,margin=2cm}
\usepackage{graphicx,array,tabularx,tabularray,colortbl,multirow,multicol,booktabs,enumitem,pifont,ulem,fancybox,diagbox,cancel}
\usepackage{pstricks,pstricks-add,pst-plot,pst-tree,pst-node,pst-text,pst-func,pst-math,pst-bspline,pst-3dplot,pst-eucl,pst-solides3d}
\let\pscalculate\undefined
\usepackage{pst-bezier}
\usepackage{tikz,pgfplots,tkz-tab}
\usetikzlibrary{arrows,arrows.meta,patterns,calc,positioning,trees,decorations.markings,shapes.geometric,shapes.misc,shapes.arrows,intersections,decorations.pathmorphing,decorations.pathreplacing,backgrounds,shadows,babel}
\pgfplotsset{compat=1.18}
\usepgfplotslibrary{fillbetween}
\usepackage[np]{numprint}
\usepackage{listings,makecell,pas-tableur,twoopt}
\let\barre\undefined
\usepackage[pstricks]{tablvar}
\definecolor{darkgreen}{RGB}{0,100,0}
\providecommand{\red}{\color{red}}
\providecommand{\blue}{\color{blue}}
\providecommand{\green}{\color{green}}
\providecommand{\black}{\color{black}}
\providecommand{\white}{\color{white}}
\providecommand{\gray}{\color{gray}}
\providecommand{\cyan}{\color{cyan}}
\providecommand{\magenta}{\color{magenta}}
\providecommand{\yellow}{\color{yellow}}
\providecommand{\violet}{\color{violet}}
\providecommand{\orange}{\color{orange}}
\providecommand{\euro}{\texteuro}
\providecommand{\EUR}{\texteuro}
\providecommand{\ieme}{\textsuperscript{e}}
\providecommand{\degres}{\ensuremath{^\circ}}
\providecommand{\degree}{\ensuremath{^\circ}}
\providecommand{\vv}[1]{\overrightarrow{#1}}
\providecommand{\Coord}[1]{#1}
\providecommand{\decosix}{}
\providecommand{\starredbullet}{\ensuremath{\bullet}}
\pagestyle{empty}
\setlength{\parindent}{0pt}
\usepackage{adjustbox}\DeclareUnicodeCharacter{2236}{:}\DeclareUnicodeCharacter{2190}{\ensuremath{\leftarrow}}
\providecommand{\R}{\mathbb{R}}
\providecommand{\N}{\mathbb{N}}
\providecommand{\Z}{\mathbb{Z}}
\providecommand{\Q}{\mathbb{Q}}
\providecommand{\D}{\mathbb{D}}
\providecommand{\C}{\mathbb{C}}
\providecommand{\e}{{\mathrm{e}}}
\providecommand{\vect}[1]{\overrightarrow{#1}}
\providecommand{\vectt}[1]{\overrightarrow{\mathrm{#1}}}
\providecommand{\vvt}[1]{\overrightarrow{\mathrm{#1}}}
\providecommand{\np}[2][]{{#2~#1}}
\providecommand{\vv}[1]{{\overrightarrow{#1}}}
\renewcommand{\d}{{\mathrm{d}}}
\providecommand{\degres}{{^\circ}}
\providecommand{\degre}{{^\circ}}
\providecommand{\blue}{\color{blue}}
\providecommand{\red}{\color{red}}
\providecommand{\enskip}{\quad}
\providecommand{\llbracket}{\lbrack\!\lbrack}
\providecommand{\rrbracket}{\rbrack\!\rbrack}
\providecommand{\euro}{{\text{€}}}
\providecommand{\no}{{\text{n°}}}
\providecommand{\interval}[2]{{[#1;#2]}}
\providecommand{\intervFO}[2]{{[#1;#2[}}
\providecommand{\dotfill}{\ldots}

\providecommand{\ds}{\displaystyle}
\providecommand{\textding}[1]{#1}
\providecommand{\Oij}{\(\mathrm{O};\vec{\imath},\vec{\jmath}\)}
\providecommand{\Oijk}{\(\mathrm{O};\vec{\imath},\vec{\jmath},\vec{k}\)}
\providecommand{\Ouv}{\(\mathrm{O};\vec{u},\vec{v}\)}
\providecommand{\intervFF}[2]{\ensuremath{[#1;#2]}}
\renewcommand{\euro}{\ifmmode\text{\texteuro}\else\texteuro\fi}
\DeclareUnicodeCharacter{00B0}{\ifmmode^{\circ}\else\textdegree\fi}
\DeclareUnicodeCharacter{20AC}{\ifmmode\text{\texteuro}\else\texteuro\fi}
% Énoncés : Éducation nationale. Transcription et corrigés : APMEP, Denis Vergès et contributeurs. Découpage et classement : CoopMaths, Sébastien Lozano, Éric Elter et contributeurs. Adaptation de la présentation : Maths Explorer.
% CC BY-SA : https://coopmaths.fr/www/about/
% https://coopmaths.fr/alea/static/bac/2026/tex/bac_2026_06_sujet2_asie_2.tex
\begin{document}
\noindent\textit{APMEP / CoopMaths. CC BY-SA. Présentation adaptée par Maths Explorer.}\par\medskip
\emph{On pourra traiter indépendamment les deux parties de l'exercice. \\
On arrondira, si nécessaire, les résultats à $10^{-3}$ près.}

\medskip
Dans cet exercice, on s'intéresse aux lancers-francs effectués par un joueur lors de compétitions de basketball.

Pour modéliser la situation, on considère dans chaque partie du problème que les conditions dans lesquelles s'effectuent ces lancers sont identiques et que ces lancers sont indépendants deux à deux.

\subsection*{Partie  A}
	Les statistiques de réussite des lancers-francs d'un joueur sont de 49,2~\% lors d'une saison.

	Dans cette partie, on assimilera cette fréquence à sa probabilité de réussite d'un lancer-franc.

	Au cours d'un match, ce joueur tente 16 lancers-francs.

	On désigne par $X$ la variable aléatoire qui donne le nombre de lancers-francs réussis par ce joueur lors de ce match.

	\begin{enumerate}
		\item Justifier que la variable aléatoire $X$ suit une loi binomiale dont on précisera les paramètres.

		\item Déterminer l'espérance de la variable aléatoire $X$ et l'interpréter dans le contexte de cet exercice.

		\item Calculer $P(X=5)$.

		\item Calculer la probabilité que le joueur réussisse au moins six lancers-francs.
	\end{enumerate}

\subsection*{Partie  B}
	On note $p$ la probabilité que le joueur réussisse un lancer-franc, où $p$ est un réel tel que
	\[0 \leqslant  p \leqslant  1\]

	On se place dans le cas où le joueur effectue 3 lancers-francs.

	On désigne par $Y$ la variable aléatoire qui donne le nombre de lancers-francs réussis par ce joueur.

	\begin{enumerate}
		\item Déterminer les valeurs prises par la variable aléatoire $Y$.

		\item Exprimer $P(Y=2)$ en fonction de $p$.

		\item Donner la loi de probabilité de $Y$.

		Présenter la réponse sous forme de tableau.

		\item Montrer que
		\[P(Y \geqslant  2)=-2 p^{3}+3 p^{2}\]

		\item On considère la fonction $f$ définie sur l'intervalle $[0 ; 1]$ par
		\[f(x)=-2 x^{3}+3 x^{2}.\]
		\begin{enumerate}
			\item Étudier le sens de variation de la fonction $f$ sur l'intervalle $[0 ; 1]$ et dresser son tableau de variation en y faisant figurer les valeurs aux bornes de l'intervalle $[0 ; 1]$.

			\item En déduire l'existence d'une unique valeur $\alpha$ dans l'intervalle $[0 ; 1]$ telle que $f(\alpha)=0,9$.

			\item Donner un encadrement d'amplitude $10^{-2}$ de cette valeur $\alpha$.

			\item Interpréter la valeur de $\alpha$ dans le contexte de l'exercice.
		\end{enumerate}
	\end{enumerate}

\end{document}
